% !PRINTED 632
\weberpage{632}
%<<CONT
the place of $\omega'$ is taken by the natural prime $p$, which is divisible
by $\mathfrak p$; and we speak accordingly of \gesperrt{the primes related
to a number $\omega$ or foreign to it}.

\gesperrt{A number $\omega$ that has no related primes at all is
a unit.}

In the field $\Omega_m$ (as indeed in every normal field) all
the conjugate numbers $\omega_n$ are related to the same primes.

\begin{center}
\S.\ 174.

\gesperrt{Kummer's Theorem.}
\end{center}

In the theory of the division of the circle (Bd.\,I, \S.\,169)
we have already become acquainted with certain decompositions of the
natural primes into factors composed of roots of unity,
and it is now of importance to investigate how these
factors stand to the prime factors in the field $\Omega_m$.

We consider a prime $p$ which belongs to the exponent 1,
so that $m$ is thus a divisor of $p-1$, and set
\[
p - 1 = m m'.
\]

From the $p$th roots of unity $\varrho$, periods of $m'$ terms each
can be formed according to Bd.\,I, \S.\,167, which, if $g$ is a primitive
root of $p$ (not, as above, of $m$), have the expression:
\weq{1}{
\begin{aligned}
\eta &= \varrho \phantom{^g} + \varrho^{g^m} + \varrho^{g^{2m}} + \cdots + \varrho^{g^{(m'-1)m}} \\
\eta_1 &= \varrho^g + \varrho^{g^{m+1}} + \varrho^{g^{2m+1}} + \cdots + \varrho^{g^{(m'-1)m+1}} \\
&\ \, \cdots\cdots\cdots\cdots\cdots\cdots\cdots\cdots\cdots\cdots \\
\eta_{m-1} &= \varrho^{g^{m-1}} + \varrho^{g^{2m-1}} + \varrho^{g^{3m-1}} + \cdots + \varrho^{g^{p-2}}.
\end{aligned}
}

If, as in \S.\,169 of the first volume, we denote by $\varepsilon$ a
primitive $(p-1)$th root of unity, so that
\weq{2}{\varepsilon^{m'} = r,}
as hitherto, is a primitive $m$th root of unity, then the
resolvents are
\begin{equation*}\tag*{3)}
(\varepsilon^\lambda, \varrho) = \varrho + \varepsilon^\lambda \varrho^g + \varepsilon^{2\lambda}\varrho^{g^2} + \cdots + \varepsilon^{(p-2)\lambda}\varrho^{g^{p-2}},
\end{equation*}
% [sic] printed "3)" without an opening parenthesis, unlike the other display numbers on this page
in which $\lambda$ may be any integral rational number, and in
particular for $\lambda = m'$:
\weq{4}{(r, \eta) = \eta + r\eta_1 + r^2\eta_2 + \cdots + r^{m-1}\eta_{m-1}.}
